\scanpage{43}\mspage{30}

\margin{\unsure{$2^{\mathrm{e}}$|$\S$}}

Démonstration du th. 2.1 et du cor. 2.5.

2.1. $1^\circ$). \deleted{On prouve d'abord le cas particulier pour $K=\varnothing$, i.e. \unsure{que|pour} $\forall$ ouvert $J$ de $\deleted{I}\unsure{I|A}$, la flèche}

\deleted{\[
 u_{\unsure{J|I}}:X_J\longrightarrow Y_J
\]}

\deleted{est dans $C$. Considérons l'ensemble \inserted{$P\subset\unsure{\mathcal P|\mathfrak P}(I)$} des ouverts $J$ de $I$ pour lesquels cette condition soit satisfaite, prouvons que \inserted{pour} toute partie $(J_\alpha)_{\alpha\in A}$ tot. ordonnée \inserted{(par inclusion)} de $P$, la réunion $\bigcup J_\alpha$ est encore dans $P$. On sait (\unsure{où|ou} on \unsure{pourra|prouve} en fait \unsure{comment}) qu'il existe, pour tt ens. tot. ordonné, \deleted{de toute partie} \inserted{une partie} bien ordonnée cofinale. Donc on peut \unsure{suppo-} $A$ bien ordonné.}

\unsure{OPS} $J=I$, \inserted{$K\subsetneqq I$.} On construit une suite transfinie \inserted{str. croissante} d'ouverts $I_\alpha$, avec

\[
 \left\{\begin{aligned}
 I_0&=\deleted{\varnothing}K\\
 I_{\alpha+1}&= I_\alpha\cup I_{/s_\alpha}\quad\text{où }s_\alpha\in I\setminus I_\alpha
 \quad(\text{si }I_\alpha\ne I,\\[-2pt]
 &&\text{sinon la \unsure{constr.}}\\[-2pt]
 &&\text{transfinie s'arrête})\\
 I_\alpha&=\bigcup_{\beta<\alpha}I_\beta\quad\text{si $\alpha$ ordinal limite.}
 \end{aligned}\right.
\]

Il existe un $\alpha_0$ tel qu'on ait $I_{\alpha_0}=I$. On voit alors comme au \S\,1

% archive sheet number: 42

